\textit{Mycobacterium tuberculosis} (M.tb) drug resistance is a major global health concern, with current genotypic drug susceptibility tools relying on near-static catalogues. These tools risk missing recently-evolved variants, leading to False Negative (FN) calls. The PE/PPE gene family comprises $\sim$10\% of the M.tb genome and is routinely masked from analysis, despite documented links to resistance and virulence \cite{marin2022benchmarking, modlin2021exact, dejesus2017comprehensive, d2023pe}. This study mined $\sim$45,000 M.tb genomes for variants enriched in FN samples, emphasising the PE/PPE gene family, and assessed a single-layer trained neural network initialised with catalogue variants. TB-Profiler \cite{phelan2019integrating} calls were compared with phenotypes to identify per-drug FN calls. A per-variant Fisher's exact test (FN vs True Negative) was conducted, controlling for significance level, association, lineage and catalogue-inclusion to produce an augmented variant set. This set was included in a single-layer network across 15 drugs, with variant ablation used to quantify each variant's independent contribution to sensitivity. The model matched or exceeded TB-Profiler sensitivity for 13--14 of 15 drugs, with the largest gains on newer drugs (Clofazimine 0.33 vs 0.166 baseline, Delamanid, Bedaquiline), at a deliberate $\sim$0.80 specificity floor. Catalogue augmentation only produced marginal aggregate change, emphasising the model's representational ceiling. Post-filtering, three variants independently contributed to sensitivity ($ethA$ (confounder), $Rv3254$, $pbpB$), while a large set of high prevalence, uncatalogued PE/PPE variants recurred across 12 drugs, mostly consistent with compensatory or epistatic roles. This approach acts as a proof of concept for a trainable, FN-informed alternative to catalogue-based genomic DST. There is a genuine novel signal concentrated in masked PE/PPE regions, motivating interaction-aware model architecture and targeted wet-lab validation.
